\section{Data and Evaluation Protocol}
\label{sec:protocol}

\subsection{Pooled document set}
Earlier versions of this study scored 77 documents and chose thresholds on the
same documents they were reported on. We replaced that setup with a pooled set
of 2{,}291 documents (1{,}143 clean, 1{,}148 poisoned) from three origins
(Table~\ref{tab:pool}). BIPIA~\cite{yi2023benchmarking} supplies clean
documents and poisoned copies with an inserted instruction; PoisonedRAG~\cite{zou2024poisonedrag}
supplies clean Natural Questions passages~\cite{kwiatkowski2019nq}, taken from
the BEIR collection~\cite{thakur2021beir}, and
poisoned passages written to steer an answer, usually without any trigger
phrase; and our own development documents supply a small in-house source, a
set of cross-domain clean documents and ten hand-written adversarial documents.
Exact duplicates were removed, and no text appears with both labels.

\begin{table}[t]
\caption{Pooled document set by origin (from \texttt{data/pooled\_stats.md}).}
\label{tab:pool}
\centering
\footnotesize
\begin{tabular}{@{}lrr@{}}
\toprule
Origin & Clean & Poisoned \\
\midrule
BIPIA & 578 & 578 \\
PoisonedRAG (NQ passages) & 498 & 500 \\
In-house development documents & 55 & 60 \\
In-house cross-domain documents & 12 & 0 \\
In-house adversarial documents & 0 & 10 \\
\midrule
Total & 1{,}143 & 1{,}148 \\
\bottomrule
\end{tabular}
\end{table}

\begin{figure*}[t]
\centering
\resizebox{0.95\textwidth}{!}{%
\begin{tikzpicture}[
  font=\footnotesize,
  >=Stealth,
  box/.style={draw=black, fill=white, rounded corners=2pt, align=center, inner sep=4pt, minimum height=0.85cm},
  role/.style={box, minimum width=4.6cm},
  arr/.style={->, thick},
  note/.style={font=\scriptsize\itshape, text=black!75, align=left}]
\node[box, minimum width=3.2cm] (pool) at (0,0) {pooled set\\[-1pt]{\scriptsize 2{,}291 documents}\\[-1pt]{\scriptsize 1{,}143 clean, 1{,}148 poisoned}};
\node[box, fill=gray!15, minimum width=3.6cm] (split) at (5.0,0) {group-aware split\\[-1pt]{\scriptsize stratified, seed 42, 43, 44}\\[-1pt]{\scriptsize a group is never divided}};
\node[role] (base) at (10.8, 3.3) {baseline\\[-1pt]{\scriptsize 285 documents}};
\node[role] (train) at (10.8, 1.1) {train\\[-1pt]{\scriptsize 1{,}116 documents}};
\node[role] (val) at (10.8,-1.1) {validation\\[-1pt]{\scriptsize 373 documents}};
\node[role, line width=1.2pt] (test) at (10.8,-3.3) {test (held out)\\[-1pt]{\scriptsize 373 documents}};
\node[box, dashed, minimum width=2.6cm] (drop) at (5.0,-3.3) {dropped\\[-1pt]{\scriptsize 144 poisoned twins}\\[-1pt]{\scriptsize of baseline documents}};
\foreach \n in {base,train,val,test}{\draw[arr] (split.east) -- ++(0.7,0) |- (\n.west);}
\draw[arr] (pool) -- (split);
\draw[arr] (split.south) -- (drop.north);
\node[note, anchor=west] at (13.3, 3.3) {reference corpus for the\\embedding signal B};
\node[note, anchor=west] at (13.3, 1.1) {fit D and the fusion; thresholds\\$\tau_\alpha$ from out-of-fold clean scores};
\node[note, anchor=west] at (13.3,-1.1) {dry runs and design checks\\(numbers are optimistic, not results)};
\node[note, anchor=west] at (13.3,-3.3) {scored once per seed after the\\freeze; the lock file refuses a rerun};
\end{tikzpicture}}
\caption{Evaluation protocol. Documents are grouped (an original and its poisoned copies stay together), split into four roles with three seeds, and every threshold is fixed on training data before the test role is scored once.}
\label{fig:protocol}
\end{figure*}

\subsection{Group-aware splits}
Documents that share an underlying text (a BIPIA original and its poisoned
copies, the paraphrases of one PoisonedRAG target) form a group, and a group is
never divided between roles (Fig.~\ref{fig:protocol}). A seeded, stratified
split assigns each group to one of four roles. The \emph{baseline} role (a
quarter of the groups that contain a clean document; 285 documents) only builds
the embedding detector's reference corpus, and any poisoned copy of a baseline
document is dropped (144 documents) so that a reference document never sits next
to its own poisoned twin. The rest is divided into \emph{train} (1{,}116
documents), \emph{validation} (373) and \emph{test} (373). We repeat the split
with three seeds (42, 43, 44). The seeds draw different group assignments from
the same pool, so the three test sets overlap heavily; agreement across seeds
shows that results are stable to the split, not that they are independent
replications. Grouped splits, a held-out role scored once and precision at low
prevalence respond to pitfalls that Arp et al.\ list for machine learning in
security: sampling bias, data snooping and the base-rate fallacy~\cite{arp2022dos}.

\subsection{Operating points and metrics}
A detector is only useful at a stated false-alarm rate, so we fix the operating
point before looking at any held-out score. For a false-alarm budget $\alpha$
(5\% and 1\%) the threshold $\tau_\alpha$ is the $(1-\alpha)$ quantile of the
\emph{training} clean-document scores, out-of-fold for every fitted method
(Eq.~\ref{eq:threshold}), and it is applied unchanged to the evaluation
documents. Fitted methods (logistic regression, the sentence classifier) are
trained on the training role only; we do not refit on training plus validation,
a conservative choice. With labels $y=1$ for poisoned and $y=0$ for clean
documents we report the recall and the realised false-alarm rate,
\begin{equation}
    \begin{gathered}
    \mathrm{TPR}=\Pr\big[s(d)>\tau_\alpha\mid y=1\big],\\
    \mathrm{FPR}=\Pr\big[s(d)>\tau_\alpha\mid y=0\big],
    \end{gathered}
    \label{eq:tprfpr}
\end{equation}
with 95\% intervals from a group bootstrap~\cite{efron1979bootstrap}: in each of $B=1{,}000$ replicates we
draw whole groups with replacement, recompute both rates, and report the 2.5th
and 97.5th percentiles,
\begin{equation}
    \mathrm{CI}_{95}=\big[\,q_{0.025}(\mathcal{T}),\ q_{0.975}(\mathcal{T})\,\big],\quad
    \mathcal{T}=\{\mathrm{TPR}^{(b)}\}_{b=1}^{B}.
    \label{eq:bootstrap}
\end{equation}

\paragraph{Precision at a given share of poisoned documents.}
Recall and false-alarm rate do not say how many alerts are real, because that
depends on how rare poisoned documents are. If a share $\pi$ of the retrieved
documents is poisoned, the precision of an alert follows from Bayes' rule,
\begin{equation}
    \mathrm{Prec}(\pi)=\frac{\pi\,\mathrm{TPR}}{\pi\,\mathrm{TPR}+(1-\pi)\,\mathrm{FPR}^{*}},
    \label{eq:precision}
\end{equation}
where $n_0$ is the number of clean evaluation documents and $\mathrm{FPR}^{*}$ is
the observed false-alarm rate, or $3/n_0$ when no false alarm was observed
(the rule-of-three upper bound~\cite{hanley1983zero}), so a method with no observed false alarm is
not flattered. The expected number of
alerts among $N$ retrieved documents is $N\,[\pi\,\mathrm{TPR}+(1-\pi)\,\mathrm{FPR}]$,
of which a share $1-\mathrm{Prec}(\pi)$ are false. We quote precision at
$\pi=1\%$ and $0.1\%$ because the alert load, not recall alone, decides whether
a detector can be deployed; Axelsson made the same point for intrusion
detection~\cite{axelsson2000base}. These values are computed from the rates above; they
are not measured on a stream with that share.

\subsection{One-shot test discipline}
The configuration (methods, budgets, seeds, bootstrap settings) was committed
and tagged before the test role was scored. The held-out documents were scored
once per seed by a single script that refuses to run without the frozen
configuration and writes a lock file recording the commit, the configuration
hash and the split hash; a second run is refused. Test results therefore
contain no choice made after seeing them. Anything we learned about a method
after the freeze is reported as a limitation, not as a fix.